% ==============================================================================
% MARCO TEÓRICO
% ==============================================================================
% Contenido provisto directamente por el estudiante (dos veces: la primera
% entrega se cortó a los 50,000 caracteres de por medio de la bibliografía;
% el reenvío completo llegó con la codificación corrupta -- "ó" en vez de
% "ó", etc. -- así que este archivo se reconstruyó a partir del texto de la
% primera entrega, que sí tenía acentos correctos, y se completó con las
% entradas bibliográficas del reenvío).
%
% Corrección aplicada: "la restricción de disponibilidad descrita en la
% sección 4.1" (en Inferencia mediante ONNX) era una referencia numérica
% hardcodeada a la subsección "Alcance técnico" de h-alcance.tex. Con
% Alcance como capítulo propio, esa subsección no será "4.1" (mismo tipo de
% problema que "la sección 1.3" en Alcance). Se agregó
% \label{subsec:alcance-tecnico} en h-alcance.tex y aquí se usa \ref.
%
% Verificación de fuentes (2026-08-15): CVE-2021-27185 confirmado por Snyk
% (command injection, paquete npm samba-client) -- coincide con el texto;
% no se pudo verificar independientemente el puntaje CVSS 9.8 exacto.
% republicamastercard2025 (200% de incremento; tecnológico 7.9%, tercer
% sector más atacado) confirmado literal. gem2024 (segunda tasa de
% actividad emprendedora temprana del mundo, 32.4%) NO se pudo verificar
% (la URL del PDF responde 403); revisar manualmente antes de la entrega.
%
% NOTA: la entrada ecosistemastartup2025 llegó en el .bib del estudiante
% pero no se cita en ningún punto de este capítulo; se agregó al .bib de
% todas formas (no genera error, biblatex solo imprime lo citado) pero
% queda sin usar hasta que se referencie en algún capítulo.
%
% Actualización 2026-09-04: el estudiante volvió a pegar el envío original
% completo (misma codificación corrupta "ó" documentada arriba). Se revisó
% contra el archivo ya reconstruido y NO se reincorporó la mayoría del
% contenido: la cifra "Node.js >40% de uso entre desarrolladores" no tiene
% cita ni en la versión original; las fuentes de blog personal (w3schools,
% Medium: cocca2025/ejaz2025/rusin2026/kalita2024, browserstack2025,
% codeb2024, socket2024, mdnwebdocs) se descartaron por calidad frente a las
% fuentes ya usadas (OpenJS Foundation, expressjs); checkpoint2023 no tiene
% URL verificable; la sección "Entorno de Prueba y Reproducibilidad"
% (Docker/Artillery/DotDotPwn/Commix) es diseño experimental y se dejó
% pendiente para el capítulo de Metodología en vez de duplicarla aquí.
%
% Se intentó agregar una sección "Trabajos previos y posicionamiento" con
% cuadro comparativo de trabajos ML previos, pero se detectó que
% e-antecedentes.tex YA tiene una sección equivalente ("Trabajos académicos
% en detección de ataques web mediante aprendizaje automático",
% \label{cuadro:trabajosprevios}) con los mismos cinco trabajos. Por
% instrucción explícita del estudiante, ese contenido se queda únicamente
% en Antecedentes y no se repite aquí.
%
% Por la misma lógica, se eliminó la subsección "RASP comerciales
% existentes" (Cuadro cuadro:rasp-comercial) de este archivo: duplicaba
% casi por completo la sección "Soluciones RASP existentes: antecedentes
% tecnológicos" (Cuadro cuadro:rasp, ahora \label{sec:rasp-antecedentes})
% de Antecedentes. Se reemplazó por una oración que remite a esa
% comparación en vez de repetirla, y se agregó el label faltante en
% e-antecedentes.tex para poder referenciarla. También se dejó fuera de
% este archivo la sección "Contexto Tecnológico y de Ciberseguridad en
% Guatemala" completa del envío del estudiante (GEM, Mastercard, MINECO,
% BID-OEA): esas mismas cifras y fuentes ya se citaron en el capítulo de
% Justificación; este archivo mantiene solo la versión condensada que ya
% tenía, que remite a Antecedentes/Justificación para el detalle.
%
% Contenido nuevo sí incorporado del envío del estudiante (2026-09-04),
% verificado antes de escribirse: cita faltante de microsoft2009stride en
% STRIDE; mención de Helmet/CORS/RateLimit/JWT/bcrypt en la subsección de
% middleware; subsección nueva "Defensa en profundidad" con figura TikZ
% (fuente cloudflaredefenseindepth NO verificable por WebFetch, 403,
% marcado en el .bib); cuarta limitación de RASP (complejidad de
% actualización del modelo embebido); indicadores de log adicionales para
% SQLi/XSS/Path Traversal/Command Injection (mismas fuentes ya verificadas,
% sin citas nuevas); referencia cruzada OWASP API3:2023 -> Path
% Traversal/LFI.
%
% Explícitamente NO incorporado por falta de fuente o por baja calidad de
% fuente: la cifra "Node.js >40% de uso entre desarrolladores" (sin cita ni
% en el envío original); las estadísticas de escala de npm de socket2024
% (2.5M paquetes, 184 mil millones de descargas/mes, 165,486 mantenedores:
% no se pudo confirmar el contenido específico por WebFetch); las cifras de
% Akamai (311/150 mil millones de ataques, +33%) vía encryptionconsulting2025
% (URL devuelve 404) y del Enterprise Strategy Group (62%) vía equixly2024
% (la portada del blog no confirma la cifra); checkpoint2023 (sin URL) y
% checkpointresearch2024 (la portada no muestra el contenido citado); y las
% fuentes de blog personal (w3schools, Medium: cocca2025/ejaz2025/
% rusin2026/kalita2024, browserstack2025, codeb2024, mdnwebdocs).
%
% Corrección 2026-10-02 (clasificación OWASP/MITRE, verificada por WebFetch
% contra las páginas oficiales antes de escribirse):
%  (1) Path Traversal (CWE-22/CWE-23) está mapeado en A01:2021 Broken Access
%      Control, NO en A03:2021 Injection -- confirmado contra
%      top10.owasp.org/2021/A01_2021-Broken_Access_Control. SQLi (CWE-89),
%      XSS (CWE-79) y OS Command Injection (CWE-78) sí están en A03:2021 --
%      confirmado contra top10.owasp.org/2021/A03_2021-Injection, que
%      además confirma que CWE-22 NO aparece en esa página. Se corrigieron
%      las subsecciones de cada vector, la sección "Marcos de referencia", y
%      se agregó el Cuadro cuadro:owasp-mitre-cwe.
%  (2) Se eliminó la afirmación (agregada por error en la sesión anterior,
%      2026-09-04) de que API3:2023 (Broken Object Property Level
%      Authorization) se corresponde con Path Traversal/LFI: verificado
%      contra api-security.owasp.org que API3:2023 cubre exposición
%      excesiva de propiedades y Mass Assignment, sin relación con
%      manipulación de rutas de archivo.
%  (3) MITRE: verificado contra attack.mitre.org que T1189 (Drive-by
%      Compromise) se centra explícitamente en comprometer el endpoint
%      cliente, con XSS como mecanismo de entrega -- se mantiene solo para
%      XSS. T1190 (Exploit Public-Facing Application) menciona
%      explícitamente SQL injection, path traversal y el OWASP Top 10/CWE
%      Top 25 -- se usa para SQLi, Path Traversal y Command Injection.
%  (4) Las consecuencias de SQLi se ajustaron a lo que A03:2021 documenta
%      (modificación/borrado de datos, invocación de procedimientos
%      almacenados) en vez de la lista genérica anterior.
%  Entradas .bib nuevas, cada una verificada por WebFetch antes de
%  escribirse: owaspa01_2021, owaspa03_2021, owaspapi3_2023, owasptraversal,
%  mitret1189, mitret1190.
%
% Propagación: el mismo error (los cuatro vectores bajo "A03:2021") estaba
% también en g-objetivos.tex (Objetivo específico 1); se corrigió ahí
% también por ser la misma afirmación factual incorrecta, no una decisión
% de alcance distinta.
% Actualización 2026-10-04: las figuras fig:defensa-profundidad y
% fig:rasp-vs-perimetral (antes TikZ) se reemplazaron por las imágenes
% provistas por el estudiante (figuras/figura2.png y figuras/figura3.png).
% El párrafo de defensa en profundidad se ajustó a las cinco capas de la
% figura nueva, y la leyenda de la figura 3 pasó de "arriba/abajo" a
% "izquierda/derecha".
\ifdefined\modoensamble
\else
	\input{0-preambulo_standalone}
\fi

\chapter{Marco teórico}

\section{Node.js y arquitectura orientada a eventos}

\subsection{Modelo de ejecución \textit{event-driven}}
Node.js es un entorno de ejecución de JavaScript, asíncrono y orientado a eventos, que ejecuta el
motor V8 de Google Chrome fuera del navegador \cite{openjsfoundation2024,nodejs2026intro}. A diferencia de
los servidores tradicionales que asignan un hilo por conexión entrante, Node.js opera con un único hilo de ejecución de JavaScript organizado en torno a un \textit{Event
Loop}, que permite realizar operaciones de entrada y salida sin bloqueo al delegarlas al kernel del
sistema siempre que es posible \cite{nodejs2026eventloop}. Según el modelo de ejecución de
JavaScript, cuando la pila de llamadas (\textit{Call Stack}) queda vacía se toma el siguiente
trabajo pendiente de la cola para ejecutarlo \cite{mdn2026execution}.

Este modelo es la razón técnica por la que Node.js es la plataforma elegida para
\textit{logSguarDian}: su capacidad de atender miles de conexiones concurrentes sin bloquear el
hilo principal permite instrumentar peticiones HTTP en tiempo real sin degradar la capacidad de
respuesta del servidor, siempre que el propio mecanismo de detección respete la naturaleza no
bloqueante del entorno. Esta restricción de diseño es la que obliga a que la inferencia del modelo
de detección se ejecute fuera del hilo principal, mediante un \texttt{worker\_thread} independiente,
tal como se detalla en la sección \ref{sec:onnx-inferencia}.

\subsection{Middleware y el pipeline de peticiones en Express}
El middleware es el mecanismo mediante el cual \textit{logSguarDian} se integra en aplicaciones
Node.js existentes. Una función middleware en Express tiene acceso al objeto de petición
(\texttt{req}), al objeto de respuesta (\texttt{res}) y a la siguiente función en la cadena; puede
modificar ambos objetos, terminar el ciclo petición-respuesta, o ceder el control a la siguiente
función mediante \texttt{next()} \cite{expressjsguide}. Cada función middleware se ejecuta en el
orden en que fue registrada, formando un pipeline en el que la petición y la respuesta atraviesan
cada capa antes de llegar al manejador de ruta final, que cierra el ciclo enviando la respuesta al
cliente.

Este mecanismo es lo que hace posible que \textit{logSguarDian} se instale como una dependencia
adicional del proyecto sin requerir modificaciones a la infraestructura existente: al registrarse
como middleware, el sistema obtiene acceso al contenido completo de cada petición antes de que
esta llegue a la lógica de negocio de la aplicación, la posición óptima para interceptar payloads
maliciosos. La Figura \ref{fig:pipeline-middleware} ilustra este pipeline y la posición que ocupa
\textit{logSguarDian} dentro de él.

\begin{figure}[H]
\centering
\begin{tikzpicture}[
  node distance=0.5cm,
  box/.style={draw, rectangle, rounded corners, minimum height=1cm, minimum width=2.6cm, align=center, font=\small},
  detect/.style={draw, rectangle, rounded corners, minimum height=1cm, minimum width=2.6cm, align=center, font=\small, fill=gray!15, thick},
  arr/.style={-{Stealth[length=2mm]}, thick}
]
\node[box] (req) {Petición HTTP\\entrante};
\node[box, right=of req] (helmet) {Middleware\\Helmet / CORS};
\node[detect, right=of helmet] (lsd) {\textit{logSguarDian}\\(captura + inferencia)};
\node[box, right=of lsd] (auth) {Middleware\\de autenticación};
\node[box, right=of auth] (route) {Route\\Handler};
\node[box, below=1.3cm of lsd, xshift=1.5cm] (res) {Respuesta al cliente};

\draw[arr] (req) -- (helmet);
\draw[arr] (helmet) -- (lsd);
\draw[arr] (lsd) -- (auth);
\draw[arr] (auth) -- (route);
\draw[arr] (route.south) -- ++(0,-0.4) -| (res.north);
\draw[arr, dashed] (lsd.south) -- ++(0,-0.6) -| node[below, font=\scriptsize, pos=0.25] {Bloqueo HTTP 403\\si se detecta amenaza} (res.north);
\end{tikzpicture}
\caption[Pipeline de middleware en Express.js]{Pipeline de middleware en Express.js y posición de \textit{logSguarDian} dentro del ciclo de vida de una petición}
\label{fig:pipeline-middleware}
\end{figure}
\vspace{-0.3cm}
{\footnotesize Nota. Elaboración propia, con base en la especificación de middleware de Express \cite{expressjsguide}.}

El ecosistema npm ofrece middleware especializado que se integra en este mismo pipeline para
reforzar la seguridad a nivel de aplicación: Helmet.js configura cabeceras HTTP de seguridad,
\texttt{cors} gestiona las políticas de origen cruzado, y \texttt{express-rate-limit} \cite{expressratelimit2026} controla la frecuencia de peticiones para
mitigar el abuso de la API, mientras que \texttt{jsonwebtoken} \cite{auth02026jsonwebtoken} y
\texttt{bcrypt} \cite{kelektiv2026bcrypt} cubren la gestión de tokens de autenticación y el hashing
de contraseñas, respectivamente. \textit{logSguarDian} se posiciona como una capa adicional
dentro de este mismo pipeline, complementaria a estos controles y no sustituta de ellos: mientras
Helmet y CORS operan sobre cabeceras y políticas de origen, y RateLimit sobre el volumen de
peticiones, \textit{logSguarDian} analiza el contenido semántico del cuerpo, los parámetros y el
comportamiento de cada petición para detectar los cuatro vectores de ataque intra-request
resumidos en el Cuadro~\ref{cuadro:vectores}.

\section{Ecosistema npm}
npm (Node Package Manager) es el gestor de paquetes oficial de Node.js: una herramienta de línea
de comandos que resuelve dependencias contra un registro centralizado de paquetes de código
abierto \cite{npmregistry2024}. La escala de este ecosistema es directamente relevante para el
proyecto en dos sentidos. Primero, define el canal de distribución de \textit{logSguarDian}: al
publicarse como paquete npm, la herramienta hereda un mecanismo de instalación estándar
(\texttt{npm install}) ya familiar para cualquier equipo que desarrolle sobre Node.js, eliminando la
fricción de adopción que sí presentan las soluciones RASP comerciales descritas en la
Sección~\ref{sec:rasp-existentes} de Antecedentes. Segundo, la resolución de dependencias mediante \texttt{package.json}
garantiza que las versiones de las herramientas empleadas en el proyecto (Express, Jest,
onnxruntime-node) queden fijadas y sean reproducibles entre entornos, condición necesaria para la
factibilidad técnica documentada en la metodología de este trabajo.

\section{Contexto de ciberseguridad en Guatemala}
\label{sec:contexto-guatemala}
El capítulo de Antecedentes de este informe documenta en detalle el panorama de amenazas
cibernéticas en Guatemala y el estado del marco regulatorio pendiente (Iniciativa de Ley 6347).
Esta sección retoma únicamente los elementos de ese contexto que son insumo directo para
decisiones de diseño de \textit{logSguarDian}.

La estructura del tejido empresarial guatemalteco es el dato que delimita el perfil de usuario
objetivo de la herramienta. Según datos del Ministerio de Economía, las micro, pequeñas y
medianas empresas representan entre el 96~\% y el 99~\% de las empresas del país, según la fuente y el año consultados
\cite{mineco2015mipyme,agn2025mipymes}, un segmento caracterizado por equipos de desarrollo pequeños y
presupuestos que rara vez contemplan una partida específica para ciberseguridad. El Reporte
Nacional GEM Guatemala 2023-2024 documenta además que el país posee la segunda tasa de
actividad emprendedora temprana más alta del mundo (32.4~\%) \cite{gem2024}, y que iniciativas
como EmprendamosGuate, coordinada por el Ministerio de Economía y PRONACOM, buscan identificar y acompañar a emprendedores con potencial de inversión \cite{pronacom}, un
ecosistema de emprendimiento por oportunidad caracterizado por el BID \cite{iadb2024}; este trabajo
asume, como supuesto de diseño y no como dato verificado, que una parte sustancial de las
aplicaciones de ese ecosistema se construye como APIs REST sobre Node.js.

Frente a este crecimiento, la exposición a amenazas ha aumentado de forma paralela. El Cyber
Insights Report Guatemala de Mastercard documenta un incremento del 200~\% en ciberataques
durante 2024 respecto al año anterior, con el sector tecnológico en tercer lugar (7.9~\%) entre los sectores más afectados, después del
financiero (19.3~\%) y el industrial (16.8~\%) \cite{mastercard2025,republicamastercard2025}. Esta combinación, alta tasa de creación de startups sobre Node.js
y crecimiento sostenido de la exposición a ataques, sin que exista aún un marco legal vigente que
incentive la inversión en controles mínimos \cite{lahora2026estancada}, es la que ancla la
justificación técnica de este proyecto en una realidad operativa concreta y no en una abstracción.

En cuanto a la madurez institucional en ciberseguridad, el informe conjunto del Banco
Interamericano de Desarrollo y la Organización de los Estados Americanos, elaborado con el
Global Cyber Security Capacity Centre de la Universidad de Oxford, sitúa a la región en un proceso
de mejora sostenida entre 2020 y 2025 en las cinco dimensiones del Modelo de Madurez de
Capacidad en Ciberseguridad (CMM), aunque persisten brechas estructurales en recursos, talento y
coordinación intersectorial \cite{bidoea2025}. Estas brechas de talento especializado son
precisamente la condición estructural que impide a las startups guatemaltecas contratar personal
dedicado a la seguridad informática, y son la que fundamenta el requisito de diseño de
\textit{logSguarDian} de operar de forma autónoma, sin necesidad de un analista que interprete
alertas ni ajuste reglas de forma continua.

\section{Fundamentos de seguridad en aplicaciones web}

\subsection{La tríada CIA}
La seguridad de la información se articula convencionalmente en torno a tres propiedades:
confidencialidad, integridad y disponibilidad, preservar las restricciones autorizadas sobre el acceso y la divulgación de la información, mantener
su exactitud y prevenir su modificación no autorizada, y
garantizar el acceso oportuno a los recursos por parte de usuarios legítimos, respectivamente \cite{nistinfosec2026}. En el contexto de una API REST construida sobre Node.js, cada uno de los cuatro
vectores de ataque cubiertos por \textit{logSguarDian} compromete predominantemente una de
estas propiedades: la inyección SQL y el Path Traversal/LFI comprometen la confidencialidad al
exponer datos no autorizados; el Cross-Site Scripting compromete la integridad al permitir la
manipulación de la sesión del usuario; y la inyección de comandos puede comprometer las tres
simultáneamente al otorgar ejecución arbitraria con los privilegios del proceso del servidor. La
disponibilidad, en particular, es una propiedad estructuralmente frágil en Node.js: al operar sobre
un único hilo de ejecución, una sobrecarga deliberada del Event Loop puede degradar o
interrumpir el servicio para la totalidad de los clientes concurrentes, razón por la cual el propio
diseño de \textit{logSguarDian} debe evitar convertirse él mismo en un vector de degradación de
disponibilidad, de ahí los criterios de latencia máxima definidos en el Alcance de este trabajo.

\subsection{Superficie de ataque en APIs REST}
La superficie de ataque de una aplicación web es el conjunto de puntos de entrada por los que un
agente malicioso puede introducir datos o extraer información del sistema. En una API REST, esta
superficie corresponde a cada endpoint HTTP expuesto junto con los métodos, parámetros,
cabeceras y cuerpos de solicitud que acepta. Las APIs presentan una superficie de ataque
estructuralmente mayor que las aplicaciones web tradicionales orientadas a la renderización de
páginas, porque cada endpoint es un punto de entrada independiente con su propio contrato de
datos, y porque la comunicación sin estado de REST implica que cada petición debe transportar
toda la información necesaria para su procesamiento, incluyendo credenciales y parámetros que un
atacante puede manipular directamente sin necesidad de interactuar con una interfaz gráfica
intermedia. Esta característica es la que justifica que el punto de instrumentación de
\textit{logSguarDian} sea el objeto \texttt{req} de cada petición entrante, y no el tráfico de red
agregado, dado que es en ese objeto donde se materializa la totalidad de la superficie de ataque
relevante para los cuatro vectores cubiertos.

\subsection{Defensa en profundidad}
La defensa en profundidad es una estrategia de ciberseguridad que implementa múltiples capas de
controles a lo largo de la infraestructura de un sistema, bajo la premisa de que ningún control
individual es infalible: ante la eventual evasión de una capa, las capas subsiguientes deben ser
capaces de detectar, contener o mitigar el avance del ataque \cite{cloudflaredefenseindepth}.
Aplicada al contexto de una API REST en Node.js, la estrategia organiza los controles en cinco
capas funcionales: en la capa de red operan los firewalls de capas 3 y 4 y los IDS de red; en la
capa perimetral, los WAF basados en firmas; en la capa de aplicación, los middlewares de
validación de entrada (Helmet, CORS) y los mecanismos de autenticación; en la capa de detección en
tiempo de ejecución, las herramientas RASP que inspeccionan la petición dentro del propio proceso;
y en la capa de datos, el cifrado en tránsito y en reposo y el control de acceso a la base de
datos. La Figura \ref{fig:defensa-profundidad} ubica a \textit{logSguarDian} dentro de este
modelo, en la capa de detección en tiempo de ejecución.

\begin{figure}[H]
\centering
\includegraphics[width=0.85\textwidth]{figura2.png}
\caption[Capas del modelo de defensa en profundidad]{Capas del modelo de defensa en profundidad aplicadas a una API REST en Node.js, con \textit{logSguarDian} en la capa de detección en tiempo de ejecución (RASP)}
\label{fig:defensa-profundidad}
\end{figure}
\vspace{-0.3cm}
{\footnotesize Nota. Elaboración propia, con base en Cloudflare \cite{cloudflaredefenseindepth}.}

Este modelo posiciona con precisión el aporte de \textit{logSguarDian}: no reemplaza los controles
de red ni los de aplicación ya existentes, sino que complementa la capa de detección y respuesta,
que en el perfil de startup guatemalteca sin equipo de seguridad dedicado descrito en el
capítulo de Antecedentes suele estar ausente por completo.

\subsection{Modelo de amenazas STRIDE aplicado al proyecto}
STRIDE es un modelo de clasificación de amenazas que organiza los riesgos de seguridad en seis
categorías: \textit{Spoofing} (suplantación de identidad), \textit{Tampering} (manipulación de
datos), \textit{Repudiation} (negación de una acción realizada), \textit{Information Disclosure}
(divulgación de información), \textit{Denial of Service} (denegación de servicio) y
\textit{Elevation of Privilege} (escalada de privilegios) \cite{microsoft2025stride}. El Cuadro \ref{cuadro:stride} mapea los
cuatro vectores de ataque cubiertos por \textit{logSguarDian} contra las categorías STRIDE que
comprometen directamente.

\begin{table}[H]
\centering
\small
\begin{tabular}{|L{3.4cm}|L{5.0cm}|L{4.6cm}|}
\hline
\textbf{Vector de ataque} & \textbf{Categoría STRIDE comprometida} & \textbf{Propiedad de la tríada CIA} \\ \hline
SQL Injection & Information Disclosure, Tampering & Confidencialidad, Integridad \\ \hline
Cross-Site Scripting & Spoofing, Information Disclosure & Confidencialidad, Integridad \\ \hline
Path Traversal / LFI & Information Disclosure & Confidencialidad \\ \hline
Command Injection & Tampering, Elevation of Privilege & Confidencialidad, Integridad, Disponibilidad \\ \hline
\end{tabular}
\caption[Mapeo de vectores de ataque a categorías STRIDE]{Mapeo de los cuatro vectores de ataque cubiertos por \textit{logSguarDian} a las categorías del modelo STRIDE}
\label{cuadro:stride}
\end{table}
{\footnotesize Nota. Elaboración propia con base en las categorías de STRIDE \cite{microsoft2025stride} y en la tríada de confidencialidad, integridad y disponibilidad \cite{nistinfosec2026}; la asignación de cada vector es criterio del autor.}

Dos categorías de STRIDE quedan explícitamente fuera de la cobertura de la primera versión de
\textit{logSguarDian}: \textit{Repudiation}, que requeriría un sistema de auditoría e integridad de
logs independiente del mecanismo de detección, y la cobertura completa de \textit{Elevation of
Privilege} más allá de la vía de Command Injection, que abarcaría vulnerabilidades de control de
acceso como BOLA (\textit{Broken Object Level Authorization}). Esta delimitación es consistente
con el Alcance definido en este informe y se documenta aquí como una decisión metodológica
explícita, no como una omisión no reconocida.

\section{Runtime Application Self-Protection (RASP)}
\label{sec:rasp}
RASP designa una tecnología de seguridad que se integra directamente dentro de la aplicación o su
entorno de ejecución, en lugar de operar como un control perimetral externo, habilitando la
detección y prevención de ataques en tiempo real desde el interior del propio proceso \cite{contrastsecurity,bermejohiguera2022}. Esta distinción arquitectónica es la que define la categoría tecnológica en la
que se posiciona \textit{logSguarDian}: a diferencia de un WAF o un IDS de red, que inspeccionan
el tráfico desde fuera del proceso y pierden visibilidad sobre el contenido cifrado o el contexto de
ejecución de la aplicación, un sistema RASP intercepta la petición dentro del propio proceso
Node.js, con acceso directo al objeto \texttt{req} antes de que cualquier capa de cifrado externa
sea relevante. La Figura \ref{fig:rasp-vs-perimetral} ilustra esta diferencia.

\begin{figure}[H]
\centering
\includegraphics[width=0.95\textwidth]{figura3.png}
\caption[Comparación arquitectónica: control perimetral vs. RASP]{Comparación arquitectónica entre un control perimetral (WAF/IDS, izquierda) y un sistema RASP operando dentro del proceso (\textit{logSguarDian}, derecha)}
\label{fig:rasp-vs-perimetral}
\end{figure}
\vspace{-0.3cm}
{\footnotesize Nota. Elaboración propia, con base en la definición de RASP de Contrast Security \cite{contrastsecurity} y el análisis
comparativo de herramientas de protección en tiempo de ejecución de Sureda Riera et al.
\cite{bermejohiguera2022}.}

La arquitectura RASP no está exenta de limitaciones que deben reconocerse en el diseño de
\textit{logSguarDian}. La primera es el \textit{overhead} de rendimiento introducido por la
instrumentación: al interceptar cada petición para analizarla antes de que llegue al route handler,
el middleware añade latencia medible. Este riesgo se mitiga mediante el procesamiento asíncrono
no bloqueante descrito en la sección \ref{sec:onnx-inferencia}, y su magnitud es precisamente uno
de los criterios de aceptación cuantitativos del proyecto. La segunda es el riesgo de falsos
positivos: un modelo que clasifica tráfico legítimo como malicioso degrada el servicio desde la
perspectiva del negocio. En \textit{logSguarDian} este riesgo depende del umbral de bloqueo del
Random Forest (\texttt{RF\_THRESHOLD}), único componente con autoridad de bloqueo; el Isolation
Forest solo marca anomalías en el registro y no bloquea peticiones
(Sección~\ref{subsec:met-f2}). La tercera es la
dependencia de cobertura: un sistema RASP únicamente detecta ataques que atraviesan el
middleware instrumentado, por lo que ataques dirigidos a la infraestructura subyacente (la
instancia de base de datos, el proveedor de nube, otras dependencias npm) quedan fuera de su
visibilidad, lo que posiciona a \textit{logSguarDian} como una capa complementaria dentro de un
modelo de defensa en profundidad (Figura~\ref{fig:defensa-profundidad}) y no como sustituto de
los controles perimetrales descritos en la Sección~\ref{sec:soluciones-existentes}. La cuarta es la
complejidad de actualización del
modelo embebido: una vez instalada la librería en un proyecto Node.js, el modelo de ML que opera
dentro de ella corresponde a la versión con la que fue entrenado y publicado, de modo que, si
emergen nuevos vectores de ataque para los que el modelo no fue entrenado, el desarrollador debe
actualizar la librería para incorporar un modelo reentrenado. A diferencia de los WAF administrados
en la nube, que actualizan sus reglas de forma centralizada y transparente, \textit{logSguarDian}
depende del ciclo estándar de actualización de paquetes npm, lo que introduce una ventana de
exposición ante amenazas emergentes entre versiones publicadas.

\section{Vectores de ataque cubiertos por logSguarDian}
El Cuadro \ref{cuadro:vectores} resume los cuatro vectores de ataque intra-request que
\textit{logSguarDian} está diseñado para detectar, su punto de materialización específico en el
ecosistema Node.js/Express y el tipo de detección que aplica el sistema.

\begin{table}[H]
\centering
\small
\begin{tabular}{|L{2.1cm}|L{3.4cm}|L{3.8cm}|L{2.1cm}|}
\hline
\textbf{Vector} & \textbf{Punto de materialización en Node.js} & \textbf{Indicador principal en logs} & \textbf{Tipo de detección} \\ \hline
SQL Injection & Parámetros de URL, cuerpo POST, query strings concatenados en consultas & Tokens SQL en parámetros; respuestas HTTP 500; longitud anómala de campo & Supervisado (clasificación) \\ \hline
Cross-Site Scripting & Campos de texto, cabeceras y parámetros renderizados sin sanitizar & Etiquetas HTML/\texttt{<script>}; codificaciones URL/Base64 en campos de texto & Supervisado (clasificación) \\ \hline
Path Traversal / LFI & Parámetros de ruta o nombre de archivo usados en el módulo \texttt{fs} & Secuencias \texttt{../}; rutas absolutas a archivos sensibles del sistema & Supervisado (clasificación) \\ \hline
Command Injection & Módulo \texttt{child\_process} (\texttt{exec}, \texttt{execSync}) con entrada no sanitizada & Metacaracteres de shell (\texttt{;}, \texttt{\textbar}, \texttt{\&\&}); nombres de comandos del SO & Supervisado (clasificación) \\ \hline
\end{tabular}
\caption[Vectores de ataque cubiertos por logSguarDian]{Vectores de ataque cubiertos por \textit{logSguarDian}: punto de materialización, indicadores en logs y tipo de detección}
\label{cuadro:vectores}
\end{table}
{\footnotesize Nota. Elaboración propia con base en OWASP \cite{owaspa01_2021,owaspa03_2021}. El Random Forest es el único componente que bloquea; el Isolation Forest marca anomalías en el registro sobre cualquier petición, sin bloquearla.}

Los indicadores del Cuadro~\ref{cuadro:vectores} que dependen de la respuesta del servidor, como
los códigos HTTP 500 o 200 o la concentración de errores en un \textit{endpoint}, son útiles para
el análisis posterior de los registros, pero no forman parte del vector de características.
\textit{logSguarDian} decide antes de que exista una respuesta, por lo que el código de estado se
excluye de la entrada de ambos modelos (Sección~\ref{subsec:met-f2}); la detección se basa
únicamente en el contenido y la estructura de la petición.

\subsection{SQL Injection}
La inyección SQL ocurre cuando datos de entrada no sanitizados se concatenan directamente en
una consulta a base de datos, permitiendo a un atacante alterar la lógica de la consulta original.
OWASP clasifica este vector (\textit{Common Weakness Enumeration}, CWE-89) dentro de la categoría A03:2021 -- Injection, y documenta
que una consulta alterada puede devolver la totalidad de los registros de una tabla y, en
escenarios más graves, modificar o eliminar datos e incluso invocar procedimientos almacenados
\cite{owaspa03_2021}. En el contexto de Node.js, el vector se materializa en cualquier endpoint
que construya consultas mediante concatenación de cadenas en lugar de consultas parametrizadas;
un atacante que envía \texttt{' OR 1=1 --} en un campo de autenticación puede lograr que la
consulta resultante devuelva la totalidad de una tabla, evadiendo el control de acceso.

Los intentos de inyección SQL dejan una huella característica en los logs de Node.js que permite
su detección mediante análisis de comportamiento: la presencia de tokens SQL en los parámetros
(comillas simples, doble guion, punto y coma, y palabras reservadas como UNION, SELECT, DROP
u OR); un patrón de respuestas HTTP 500 concentradas en un mismo endpoint en un período corto,
generado cuando el payload malformado llega a la capa de consulta sin ser filtrado; y una longitud
de campo anómalamente mayor a la esperada, por ejemplo un campo de usuario que normalmente
recibe entre 5 y 30 caracteres y recibe un payload de 200 o más durante un ataque. La combinación
de estos tres indicadores es la firma comportamental que \textit{logSguarDian} extrae para
clasificar este vector mediante el modelo supervisado.

\subsection{Cross-Site Scripting (XSS)}
El XSS consiste en la inserción de código JavaScript malicioso en contenido que posteriormente es
renderizado por otros usuarios, en sus variantes reflejada, almacenada y basada en DOM. OWASP
lo clasifica (CWE-79) dentro de A03:2021 -- Injection \cite{owaspa03_2021}. Dentro de la
taxonomía MITRE ATT\&CK, este vector se corresponde con la técnica T1189
(\textit{Drive-by Compromise}), cuya documentación oficial identifica explícitamente la inserción
de scripts maliciosos mediante interfaces de aplicaciones web que permiten contenido controlado
por el usuario como uno de los mecanismos de entrega de esta técnica \cite{mitret1189}. Esta
técnica se centra en comprometer el equipo cliente del usuario que visita el sitio, no el servidor,
lo que es coherente con el hecho de que el XSS afecta primariamente a los usuarios de la
aplicación y no a la infraestructura que la aloja. En una API REST, el riesgo se materializa
cuando la API sirve contenido dinámico renderizado directamente en el navegador del cliente; un atacante que inyecta un script
almacenado puede robar tokens de sesión o \textit{JSON Web Tokens} (JWT) y ejecutar acciones autenticadas en nombre del
usuario comprometido.

En los logs de Node.js, los intentos de XSS presentan tres indicadores observables. El primero es
la presencia de etiquetas o atributos HTML en campos que por diseño de la API solo deberían
recibir texto plano, como \texttt{<script>}, \texttt{onerror=}, \texttt{onload=} o \texttt{javascript:}.
El segundo es el uso de codificaciones para evadir filtros superficiales, como codificación URL
(\texttt{\%3Cscript\%3E}) o de entidades HTML (\texttt{\&lt;script\&gt;}). El tercero es un patrón
temporal doble característico del XSS almacenado: una petición POST o PUT con el payload hacia
un endpoint de escritura, seguida de peticiones GET al recurso donde ese contenido queda
almacenado.

\subsection{Path Traversal / Local File Inclusion (LFI)}
El Path Traversal consiste en la manipulación de rutas de archivo mediante secuencias como
\texttt{../} para acceder a archivos fuera del directorio raíz autorizado por la aplicación \cite{owasptraversal}. LFI es un vector relacionado que ocurre cuando la aplicación
no solo lee, sino que incluye y evalúa el archivo referenciado en su propio contexto de
ejecución, lo que puede escalar a ejecución remota de código. Ambos vectores comparten la misma
causa raíz: la ausencia de validación y canonicalización de la entrada de usuario antes de
construir rutas del sistema de archivos.

A diferencia de los otros tres vectores, el OWASP Top 10:2021 no clasifica el Path Traversal
como una categoría de inyección: los identificadores CWE-22 (\textit{Path Traversal}) y CWE-23
(\textit{Relative Path Traversal}) están mapeados en A01:2021 -- Broken Access Control
\cite{owaspa01_2021}, mientras que SQL Injection (CWE-89), XSS (CWE-79) y la inyección de comandos
del sistema operativo (CWE-78) están mapeados en A03:2021 -- Injection \cite{owaspa03_2021}. La
inclusión del Path Traversal/LFI entre los vectores cubiertos por \textit{logSguarDian} no se
justifica, por lo tanto, por pertenecer a la misma categoría OWASP que los demás, sino por
compartir con ellos una causa raíz operativa (entrada controlada por el cliente que llega sin
validación a un recurso sensible del servidor) y una huella observable en el objeto \texttt{req}
que un mismo extractor de características puede capturar. En Node.js, el módulo \texttt{fs} y los
servidores de archivos estáticos son los puntos de integración más comunes donde este vector se
materializa cuando la entrada no es sanitizada.

Las técnicas de evasión documentadas añaden complejidad a la detección mediante reglas
estáticas: codificación URL simple (\texttt{\%2e\%2e\%2f} como equivalente de \texttt{../}),
codificación URL doble (\texttt{\%252e\%252e\%252f}), codificación Unicode (\texttt{..\%c0\%af}) e
inserción de bytes nulos para truncar extensiones verificadas por el servidor \cite{owasptraversal}. Esta variabilidad sintáctica es la que hace insuficiente una lista de denegación
basada en firmas fijas, y la que justifica un clasificador supervisado entrenado sobre la
distribución estadística de estos patrones. Los indicadores observables en los logs de Node.js son
tres: la secuencia \texttt{../} y sus variantes codificadas en parámetros que por diseño solo deberían
recibir nombres de archivo; la aparición de rutas absolutas o nombres de archivos sensibles del
sistema operativo (\texttt{/etc/passwd}, \texttt{.env}, \texttt{config.json}); y una respuesta HTTP 200
con contenido no esperado, o 500 cuando la ruta construida no existe.

\subsection{Command Injection}
La inyección de comandos del sistema operativo ocurre cuando una aplicación transmite entrada de
usuario no sanitizada a un intérprete de comandos del sistema; OWASP la clasifica (CWE-78) dentro
de A03:2021 -- Injection \cite{owaspa03_2021}. En Node.js, la
superficie de ataque se concentra en el módulo \texttt{child\_process}, cuyas funciones
\texttt{exec} y \texttt{execSync} invocan un shell del sistema operativo que interpreta
metacaracteres como \texttt{;}, \texttt{\textbar}, \texttt{\&\&} y \texttt{\$()}. Esta clase de
vulnerabilidad está documentada en incidentes reales del ecosistema npm: CVE-2021-27185 afectó
al paquete \texttt{samba-client}, donde el uso de \texttt{process.exec} sin sanitización de entrada
permitía inyección de comandos con un puntaje CVSS crítico de 9.8 \cite{gitlabadvisory2021}.

La variabilidad sintáctica de los payloads entre plataformas aumenta la complejidad de su
detección mediante reglas estáticas: en Linux, los metacaracteres más frecuentes son \texttt{;},
\texttt{\textbar}, \texttt{\&\&}, \texttt{\textbar\textbar} y la sustitución de comandos mediante
\texttt{`cmd`} o \texttt{\$(cmd)}; en Windows, los equivalentes son \texttt{\&}, \texttt{\textbar},
\texttt{\&\&} y \texttt{\textbar\textbar} aplicados sobre \texttt{cmd.exe}. Los indicadores observables
en los logs de Node.js son la presencia de estos metacaracteres en campos que por diseño deberían
recibir valores alfanuméricos, y la aparición de nombres de comandos del sistema operativo
(\texttt{cat}, \texttt{whoami}, \texttt{wget}, \texttt{curl}) dentro de parámetros de la API.

\section{Marcos de referencia para la selección de vectores}
La selección de los cuatro vectores cubiertos por \textit{logSguarDian} no responde a un criterio
arbitrario, sino a su correspondencia con marcos de referencia reconocidos internacionalmente y
basados en evidencia empírica de incidentes reales. Esa correspondencia no es uniforme: los
cuatro vectores no pertenecen a una única categoría de OWASP, y el Cuadro~\ref{cuadro:owasp-mitre-cwe}
documenta con precisión dónde se ubica cada uno.

El OWASP Top 10 es una lista de las diez categorías de riesgo más críticas en aplicaciones web,
construida a partir del análisis de vulnerabilidades reportadas en la industria \cite{owasp2021}.
Este trabajo emplea la edición 2021, sobre la que se definieron el protocolo y la taxonomía del
dataset; la edición 2025 ya está disponible \cite{owasp2025top10}. En ella Injection pasa de A03 a A05 y
conserva los CWE de inyección de comandos, XSS y SQL (CWE-78, CWE-79 y CWE-89), y Path Traversal
(CWE-22) sigue dentro de A01, por lo que la asignación de los cuatro vectores no cambia de fondo; la
edición 2025 no se adoptó para mantener la trazabilidad con las decisiones ya tomadas. En la edición 2021, tres de los cuatro vectores se agrupan bajo A03:2021 --
Injection: la inyección SQL (CWE-89), el XSS (CWE-79) y la inyección de comandos del sistema
operativo (CWE-78) \cite{owaspa03_2021}. El Path Traversal (CWE-22) no pertenece a esa categoría:
está mapeado en A01:2021 -- Broken Access Control \cite{owaspa01_2021}. La primera versión de
\textit{logSguarDian} cubre, por lo tanto, vectores de dos categorías OWASP distintas, y la
justificación de agruparlos en una misma herramienta es operativa (todos materializan su payload
en el objeto \texttt{req} y son detectables a partir de las mismas características extraídas), no
taxonómica.

\begin{table}[H]
\centering
\small
\begin{tabular}{|L{3.4cm}|L{2.2cm}|L{4.4cm}|L{2.4cm}|}
\hline
\textbf{Vector} & \textbf{CWE principal} & \textbf{OWASP Top 10:2021} & \textbf{MITRE ATT\&CK} \\ \hline
SQL Injection & CWE-89 & A03 -- Injection & T1190 \\ \hline
Cross-Site Scripting & CWE-79 & A03 -- Injection & T1189 \\ \hline
Path Traversal / LFI & CWE-22 & A01 -- Broken Access Control & T1190 \\ \hline
Command Injection & CWE-78 & A03 -- Injection & T1190 \\ \hline
\end{tabular}
\caption[Mapeo de vectores a CWE, OWASP Top 10 y MITRE ATT\&CK]{Mapeo de los cuatro vectores cubiertos por \textit{logSguarDian} a CWE, OWASP Top 10:2021 y MITRE ATT\&CK}
\label{cuadro:owasp-mitre-cwe}
\end{table}
\vspace{-0.3cm}
{\footnotesize Nota. Elaboración propia, con base en OWASP \cite{owaspa01_2021, owaspa03_2021} y MITRE \cite{mitret1189, mitret1190}. Dentro de las categorías del Top 10:2021 usadas aquí no existe un CWE específico de inclusión local
de archivos. A01 mapea CWE-22 \cite{owaspa01_2021} y A03 mapea CWE-98 \cite{owaspa03_2021}, que
MITRE define para la inclusión de archivos en PHP y que documenta ``local file inclusion'' solo como
término alterno \cite{mitre2026cwe98}. LFI se asocia aquí a CWE-22 \cite{mitre2026cwe22} como
variante de la manipulación de rutas.}

El OWASP API Security Top 10, en su edición 2023, reconoce que la superficie de ataque de las
APIs modernas difiere estructuralmente de las aplicaciones web tradicionales, y posiciona como
riesgo principal la Autorización Rota a Nivel de Objeto (BOLA, API1:2023) \cite{owaspapi2023};
esta categoría no está cubierta por \textit{logSguarDian} en su primera versión, consistente con
la delimitación de alcance ya señalada para el modelo STRIDE. La categoría API3:2023 (Broken
Object Property Level Authorization), por su parte, abarca la exposición de propiedades sensibles
de un objeto y la asignación masiva de propiedades que el cliente no debería poder modificar
\cite{owaspapi3_2023}; es un riesgo de autorización a nivel de propiedad y no guarda relación con
la manipulación de rutas de archivo, por lo que ninguno de los cuatro vectores cubiertos tiene un
equivalente directo en el API Security Top 10 2023.

El framework MITRE ATT\&CK es una base de conocimiento pública que documenta tácticas,
técnicas y procedimientos utilizados por adversarios reales en incidentes observados
\cite{mitreattack}. Para el dominio Enterprise, relevante para este proyecto, la táctica de Acceso
Inicial es la que agrupa los vectores cubiertos. La técnica T1190 (\textit{Exploit Public-Facing
Application}) describe la explotación de debilidades en sistemas expuestos a Internet, incluyendo
sitios web y bases de datos, y remite al OWASP Top 10 como referencia de las vulnerabilidades web
más comunes \cite{mitret1190}; es la técnica que corresponde a SQL Injection, Path Traversal y
Command Injection. La técnica T1189 (\textit{Drive-by Compromise}) se reserva para el XSS, dado
que su documentación oficial lo menciona como mecanismo de entrega y se centra en el compromiso
del cliente y no del servidor \cite{mitret1189}. La alineación de los vectores seleccionados con
ambos marcos garantiza que la cobertura de amenazas de \textit{logSguarDian} responde a evidencia
empírica de incidentes reales, y no a un criterio de selección arbitrario.

\section{Fundamentos de detección basada en aprendizaje automático}
\label{sec:ml-conceptos}
La detección de los cuatro vectores cubiertos por \textit{logSguarDian} combina dos paradigmas de
aprendizaje automático que responden a dos naturalezas distintas de amenaza. Esta sección
introduce ambos paradigmas al nivel conceptual necesario para comprender las decisiones de
arquitectura del sistema; el detalle del entrenamiento, la selección de hiperparámetros y la
evaluación comparativa de modelos candidatos se documenta en el trabajo de graduación de Diego
Valenzuela, correspondiente al componente de Inteligencia Artificial del proyecto.

\subsection{Clasificación supervisada}
El aprendizaje supervisado construye un modelo a partir de un conjunto de datos etiquetado, en el
que cada muestra de tráfico HTTP está marcada como legítima o como perteneciente a una de las
cuatro categorías de ataque cubiertas. Random Forest (RF), el algoritmo de clasificación supervisada considerado para \textit{logSguarDian}, es un método de ensamble que combina las predicciones de
múltiples árboles de decisión entrenados sobre subconjuntos aleatorios de características y
muestras, reduciendo el sobreajuste que afecta a un árbol de decisión individual
\cite{breiman2001random}. Este paradigma es apropiado para los cuatro vectores cubiertos porque cada uno
presenta patrones sintácticos y estructurales relativamente consistentes, discernibles como
características extraíbles del objeto \texttt{req} (presencia de tokens SQL, metacaracteres de
shell, secuencias de \textit{traversal}), sobre los cuales un clasificador supervisado puede
aprender fronteras de decisión efectivas.

\subsection{Detección de anomalías no supervisada}
El aprendizaje no supervisado, a diferencia del anterior, no requiere que las muestras de
entrenamiento estén etiquetadas: aprende la distribución estadística del comportamiento
considerado normal y señala como anómala cualquier observación que se desvíe significativamente de
ella. Isolation Forest (IF), el algoritmo de detección de anomalías empleado en \textit{logSguarDian},
aísla observaciones mediante particiones aleatorias recursivas del espacio de características: las
anomalías, al ser estadísticamente distintas de la mayoría de los datos, requieren en promedio
menos particiones para quedar aisladas, lo que permite asignar una puntuación de anomalía sin
etiquetas \cite{liu2008isolation}. Este paradigma complementa al clasificador supervisado, que
asigna cualquier entrada a la clase conocida más cercana aunque no se parezca a ninguna de las
vistas durante el entrenamiento. Su papel en el sistema es, sin embargo, acotado: el Isolation
Forest no bloquea peticiones, solo las marca como anómalas en el registro (\texttt{pass\_anomaly}),
y su contribución a la detección de variantes no vistas solo se midió, de forma indirecta, sobre el
conjunto de validación del componente de Inteligencia Artificial \cite{valenzuela2026ia}
(Sección~\ref{sec:met-validez}). La correspondencia entre cada paradigma y el tipo de indicador
de seguridad que produce se documenta en la Sección~\ref{sec:ioc-ioa}.

\subsection{Inferencia en tiempo de ejecución mediante ONNX}
\label{sec:onnx-inferencia}
Independientemente del algoritmo de origen, el componente de inferencia de \textit{logSguarDian}
consume un modelo serializado en formato ONNX (\textit{Open Neural Network Exchange}), un
formato estándar abierto para representar modelos de aprendizaje automático de forma
interoperable entre distintos frameworks de entrenamiento \cite{onnx2026}. La librería carga este
archivo mediante \texttt{onnxruntime-node} \cite{onnxruntimedocs} y ejecuta la inferencia dentro de
un \texttt{worker\_thread} de Node.js independiente del hilo principal; los worker threads están
pensados para operaciones de JavaScript intensivas en CPU \cite{nodejs2026worker}, y con esta
decisión de diseño se busca que la carga computacional del modelo no ocupe el Event Loop durante el
ciclo de vida de la petición. Esta decisión de arquitectura es la aplicación directa, a nivel de
implementación, de la restricción de disponibilidad descrita en la Sección~\ref{subsec:alcance-tecnico}:
sin este aislamiento, la propia herramienta de detección se convertiría en la causa de la
degradación de servicio que está diseñada para prevenir.

\section{Logs de servidor y análisis de comportamiento}
Los logs de servidor son la fuente primaria de información sobre el comportamiento de una
aplicación en producción. En Node.js y Express, los logs de acceso registran método, URL, código
de estado, tamaño de la solicitud y respuesta, dirección IP y tiempo de respuesta; los logs de error
capturan excepciones y fallos de validación que pueden revelar intentos de explotación, un volumen
inusualmente alto de respuestas 400 o 500 en un endpoint específico es indicador temprano de
sondeo con payloads malformados.

El análisis de comportamiento parte de establecer una línea base de tráfico normal, contra la cual
se evalúan desviaciones. Estos patrones pueden ser volumétricos (alta frecuencia de solicitudes
desde una misma IP, picos de errores 401/403) o semánticos (caracteres asociados a inyección en
parámetros, longitud anómala de campos, codificaciones inusuales). \textit{logSguarDian} trabaja
solo con los patrones semánticos de cada petición individual: el clasificador supervisado y el
detector de anomalías reciben ambos el vector de características de una sola petición, y las
características temporales del extractor (ventanas de 1, 5 y 60~s) no forman parte de su entrada.
Los patrones volumétricos y de secuencia, que requieren correlacionar varias peticiones, quedan
fuera del alcance de esta versión, junto con la denegación de servicio
(Sección~\ref{subsec:fuera-alcance}).

Morgan es el middleware de registro de peticiones HTTP de referencia en el ecosistema Express,
distribuido bajo la organización oficial \texttt{expressjs} en npm \cite{expressjsmorgan}. Se apoya
en el sistema de \textit{streams} de Node.js para escribir el registro sin bloquear la respuesta, y
permite definir tokens personalizados mediante \texttt{morgan.token()} para capturar campos que sus
formatos predefinidos no incluyen \cite{expressjsmorgan}. Es, por lo tanto, el punto de referencia
de cómo se registra el tráfico en una aplicación Express típica. \textit{logSguarDian} no depende
de Morgan ni extiende su registro: intercepta el objeto \texttt{req} dentro de su propia función
middleware, extrae el vector de características antes de que el manejador de ruta procese la
petición y almacena cada detección en su propio registro de eventos en SQLite. La diferencia es de
propósito: Morgan registra el tráfico para su consulta posterior, mientras que
\textit{logSguarDian} analiza cada petición para decidir, en el momento, si la bloquea.

\section{Indicadores de Compromiso (IoC) e Indicadores de Ataque (IoA)}
\label{sec:ioc-ioa}
Un Indicador de Compromiso (IoC) es un artefacto técnico u observable que sugiere que un ataque es
inminente o está en curso, o que ya pudo haberse producido un compromiso
\cite{nistioc2025,crowdstrikeioc}. Un Indicador de Ataque (IoA), en cambio, se centra en identificar
la intención y el comportamiento de un adversario mientras el ataque está en curso, enfoque que CrowdStrike presentó en 2016 como alternativa a la detección basada en IoC, a los que
critica por equivaler a firmas \cite{crowdstrike2016ioa,crowdstrikeioa}. Esta distinción es operativamente central para
\textit{logSguarDian}, que opera dentro del proceso Node.js e intercepta cada petición antes de que
llegue al manejador de ruta: las alertas que produce son, por construcción, indicadores de ataque,
porque se generan mientras el ataque está en curso y no a partir de evidencia forense posterior.
Ambos componentes producen este tipo de indicador: el Random Forest al clasificar y bloquear la
petición, y el Isolation Forest al marcarla como anómala en el registro. El
Cuadro~\ref{cuadro:ioc-ioa} sintetiza la diferencia entre ambos tipos de indicador.

\begin{table}[H]
\centering
\small
\begin{tabular}{|L{3.8cm}|L{4.4cm}|L{4.6cm}|}
\hline
\textbf{Dimensión} & \textbf{Indicador de Compromiso (IoC)} & \textbf{Indicador de Ataque (IoA)} \\ \hline
Naturaleza & Reactivo & Proactivo \\ \hline
Momento de detección & Típicamente después del incidente & Durante el ataque \\ \hline
Base de conocimiento & Firmas y artefactos conocidos & Comportamiento y patrones de la petición \\ \hline
Efectividad ante variantes desconocidas & Baja & Alta \\ \hline
\end{tabular}
\caption[Comparativa entre IoC e IoA en logSguarDian]{Comparativa entre Indicadores de Compromiso (IoC) e Indicadores de Ataque (IoA)}
\label{cuadro:ioc-ioa}
\end{table}
\vspace{-0.3cm}
{\footnotesize Nota. Elaboración propia, con base en CrowdStrike \cite{crowdstrikeioc, crowdstrikeioa}.}

\section{Soluciones de seguridad existentes}
\label{sec:soluciones-existentes}
El Cuadro \ref{cuadro:soluciones} compara las categorías de solución de seguridad disponibles
frente a \textit{logSguarDian} en las dimensiones más relevantes para el contexto de una startup
guatemalteca sin equipo de seguridad dedicado.

\begin{table}[H]
\centering
\small
\begin{tabular}{|L{2.6cm}|L{2.6cm}|L{1.6cm}|L{1.4cm}|L{1.9cm}|L{1.8cm}|}
\hline
\textbf{Solución} & \textbf{Tipo de análisis} & \textbf{Visibilidad capa 7} & \textbf{Equipo dedicado} & \textbf{Costo} & \textbf{Integración nativa en Node.js} \\ \hline
WAF (Cloudflare, AWS WAF) & Firmas estáticas & Parcial & Sí & Alto & No (perimetral) \\ \hline
IDS (Suricata, Snort) & Tráfico de red & No (ciego a HTTPS) & Sí & Medio & No (nivel de red) \\ \hline
SIEM (Splunk, Wazuh) & Correlación de eventos & Parcial & Sí (centro de operaciones de seguridad, SOC) & Muy alto & No (infraestructura externa) \\ \hline
\textit{logSguarDian} & Clasificación ML por petición & Sí (dentro del proceso) & No & Gratuito & Sí (middleware nativo) \\ \hline
\end{tabular}
\caption[Comparativa de soluciones de seguridad para APIs REST en Node.js]{Comparativa de soluciones de seguridad existentes frente a \textit{logSguarDian}}
\label{cuadro:soluciones}
\end{table}
\vspace{-0.3cm}
{\footnotesize Nota. Elaboración propia. Las calificaciones (Alto, Medio, Muy alto, equipo dedicado, visibilidad de
capa 7) son criterio del autor; el costo relativo del SIEM se apoya en ENISA \cite{enisa2021smes}.}

\subsection{Web Application Firewall (WAF)}
Un WAF opera como intermediario entre clientes y servidor, inspeccionando el tráfico HTTP/HTTPS
contra un conjunto de reglas y firmas predefinidas, principalmente en la capa 7 del modelo OSI.
Los WAF basados en firmas y en el OWASP Core Rule Set (CRS) requieren ajuste continuo y mantenimiento especializado para ser efectivos sin generar un volumen
excesivo de falsos positivos: el propio proyecto CRS advierte que los niveles de paranoia más altos
aumentan la posibilidad de bloquear tráfico legítimo por falsas alarmas y que puede ser necesario
agregar reglas de exclusión para aplicaciones con entradas complejas \cite{crs2026paranoia}. Por
construcción, una regla solo detecta lo que su patrón describe, por lo que estos WAF son poco
eficaces contra vulnerabilidades de lógica de negocio o variantes de ataque que no coinciden con
ninguna regla existente. Estas
dos limitaciones, mantenimiento especializado y ceguera ante variantes desconocidas, son
precisamente las que un modelo de aprendizaje automático entrenado sobre la distribución
estadística de los patrones de ataque, en lugar de un conjunto fijo de reglas, está diseñado para
mitigar.

\subsection{Sistemas de Detección de Intrusiones (IDS)}
Un IDS monitorea tráfico de red o actividad de sistema en busca de comportamientos sospechosos
\cite{nist80094}. Los IDS de red pierden visibilidad completa frente al tráfico cifrado mediante
HTTPS, dado que sin descifrar la conexión TLS (\textit{Transport Layer Security}) solo observan metadatos y patrones del flujo
\cite{papadogiannaki2021survey}, una ceguera crítica para detectar ataques de capa de aplicación como SQLi o XSS, donde el vector
malicioso está en el contenido del cuerpo de la solicitud. Adicionalmente, en la revisión realizada para este trabajo no se identificó ningún IDS convencional,
de red o de host, diseñado para integrarse en el modelo de concurrencia orientado a eventos de
Node.js, lo que deja una brecha arquitectónica que es la que justifica el enfoque de
\textit{logSguarDian} como middleware nativo: en lugar de observar el tráfico desde fuera del
proceso, opera desde dentro, con acceso pleno al contenido descifrado y al contexto de ejecución.

\subsection{Sistemas de Gestión de Eventos e Información de Seguridad (SIEM)}
Un SIEM agrega, correlaciona y analiza eventos de seguridad provenientes de múltiples fuentes de
infraestructura. Su principal limitación para el contexto de startups guatemaltecas es el costo: para las
pequeñas y medianas empresas, los SIEM representan inversiones considerables y, como muchas
soluciones de ciberseguridad, requieren conocimiento especializado para implementarse y
administrarse \cite{enisa2021smes}. Las soluciones SIEM ofrecidas como servicio implican además el
envío de los registros a una infraestructura externa, lo que introduce costo operativo y
consideraciones de privacidad de datos.

El capítulo de Antecedentes documenta en detalle la comparación entre las soluciones RASP
comerciales existentes (Contrast Protect, Datadog, OpenRASP) y \textit{logSguarDian}
(Sección~\ref{sec:rasp-existentes} y su cuadro comparativo); esta sección no la repite y remite a
esa comparación para el detalle de soporte de Node.js, costo y dependencia de infraestructura
externa de cada alternativa.

\section{Herramientas de simulación, pruebas y serialización}
La construcción del entorno de validación de \textit{logSguarDian} se apoya en un conjunto de
herramientas de código abierto seleccionadas por su reproducibilidad y su adopción estándar en el
ecosistema Node.js. El entorno de simulación se implementa mediante Docker Compose, lo que
garantiza que el entorno de generación de tráfico y ataques sea reproducible en cualquier momento
del proyecto y por cualquier integrante del equipo \cite{dockerdocs}; esta práctica de
contenerización como mecanismo de reproducibilidad computacional está documentada en la
literatura de ciencia de datos como una condición necesaria para la validez de resultados
experimentales dependientes de un entorno específico \cite{moreau2023}. La generación de tráfico
legítimo se realiza mediante Artillery, una herramienta de pruebas de carga que define escenarios
de tráfico mediante configuración declarativa \cite{artillerydocs}. Para el tráfico malicioso, el
protocolo contemplaba sqlmap, OWASP ZAP y Commix \cite{sqlmap, owaspzap, stasinopoulos2019}. En la
ejecución, los payloads de SQL Injection del corpus de evaluación limpio se obtuvieron con sqlmap
y los de Path Traversal y Command Injection con generadores propios, deliberadamente no derivados
de listas públicas; OWASP ZAP se dirigió mediante su API REST (Sección~\ref{sec:met-instrumentos}). Commix
interviene de forma indirecta: la lista \texttt{command-injection-commix.txt} de SecLists, derivada
de esa herramienta, es una de las fuentes de la clase \texttt{cmdi} del corpus de entrenamiento y
del corpus de la Ronda~4.

La validación de la librería distribuida como paquete npm emplea Jest como framework de pruebas
unitarias e de integración, combinado con Supertest para ejecutar peticiones HTTP contra la
aplicación sin necesidad de levantar un servidor real \cite{jestdocs, supertestdocs}. Esta
combinación permite verificar de forma determinista que el middleware intercepta peticiones con
payloads maliciosos conocidos y responde con los códigos de estado correctos ante cada categoría de
ataque. La latencia no se verifica con estas pruebas, sino con el benchmark de carga descrito en la
Sección~\ref{subsec:met-f4}.

\ifdefined\modoensamble
\else
	\section*{Referencias}
	\printbibliography[heading=none]
	\end{document}
\fi
